% =============================================================================
% Chapter 5 -- Testing and Evaluation
% Handbook guidance:
%   * EVALUATION IS NOT SELF-REFLECTION. Self-reflection lives in the
%     BCS chapter at the end. THIS chapter must ask "did the artefact
%     meet its objective requirements?" -- with measurements.
%   * Distinguish: (a) correctness tests of the implementation,
%                  (b) experimental measurements that answer research
%                      questions, (c) comparison against baselines.
% =============================================================================
\chapter{Testing and Evaluation}
\label{ch:eval}

This chapter answers two questions in order. First: is the
implementation behaving the way the methods say it should?
Second: assuming the answer to the first question is yes, what do
the experiments report about the behaviour of the methods at the
budget studied? The chapter is deliberately measurement-heavy.
Self-reflection on the project as a whole is reserved for
\Cref{ch:bcs}.

\section{Correctness Tests}
\label{sec:correctness}

Before the experimental study, the implementation was sanity-checked
against the published COOM behaviour: regularisation-loss curves were
inspected to verify that EWC's Fisher penalty fires as intended at task
boundaries (see \Cref{fig:lossreg}); training curves were spot-checked
against figures in the COOM paper for one task and seed; the
no-regularisation Fine-Tuning baseline was verified to reduce to vanilla
SAC by setting \texttt{cl\_method=None}.

\section{Quantitative Results}
\label{sec:results}

\subsection{Scope of Reproduction}
\label{sec:reproduction_scope}

The COOM benchmark paper \citep{tomilin2024coom} reports
each method's average performance across seven sequence
families, ten random seeds and a per-task budget of
$200{,}000$ steps with held-out evaluation enabled. The
present study runs each method on a single sequence (CO4),
with one to three seeds, at $20{,}000$ steps per task and
without held-out evaluation. The numbers reported in this
chapter are therefore not directly comparable to the
published headline scores in row order. Two senses of
``reproduction'' need to be kept apart:

\begin{itemize}
    \item \textbf{Pipeline reproduction} --- ``does the
        published code, with the published hyperparameters,
        run end-to-end on consumer hardware and produce the
        qualitative behaviour the methods are supposed to
        produce?'' \emph{Yes:} regularisation losses fire at
        task boundaries (\Cref{fig:lossreg}), the cumulative
        Fisher dynamics match the diagnosis of
        \citet{kirkpatrick2017ewc}, and the per-task success
        rate ordering between Fine-Tuning and the
        regularisation family is non-trivial. This satisfies
        requirement~\ref{req:pipeline}.
    \item \textbf{Numerical reproduction} --- ``does the
        published EWC score of $0.54$ on COOM appear in our
        runs?'' \emph{No, and not claimed.} At $20{,}000$
        steps per task the EWC last-five-epoch CO4 average
        is $0.037$, an order of magnitude below the
        published value, and the ordering against
        Fine-Tuning is inverted. Both differences are
        consistent with the smaller training budget; neither
        is evidence against the original publication.
\end{itemize}

The remainder of this chapter therefore reports the
\emph{within-study} ordering of the four methods rather
than the gap to the COOM paper's table. The
forking-paths analysis in \Cref{sec:ewc_online} and the
threats-to-validity discussion in \Cref{sec:validity}
re-state these scope limits when they bear on the
conclusions.

\subsection{Average Performance}

\begin{table}[t]
    \centering
    \caption{Last-five-epoch mean training success rate per task, averaged
        across the four CO4 tasks. EWC reports mean $\pm$ standard
        deviation across three seeds; the other methods are single-seed.}
    \label{tab:avg_perf}
    \IfFileExists{tab1_avg_perf.tex}{%
        \begin{tabular}{lcccc|c}
\toprule
Method & T0 & T1 & T2 & T3 & \textbf{Avg.\ Perf.} \\
\midrule
FT (1 seed) & 0.054 & 0.098 & 0.019 & 0.386 & \textbf{0.139} \\
L2 (1 seed) & 0.035 & 0.035 & 0.000 & 0.012 & \textbf{0.020} \\
MAS (1 seed) & 0.057 & 0.092 & 0.000 & 0.010 & \textbf{0.040} \\
EWC (3 seeds) & 0.028 $\pm$ 0.012 & 0.056 $\pm$ 0.005 & 0.000 $\pm$ 0.000 & 0.063 $\pm$ 0.033 & \textbf{0.037} \\
EWC-Online (1 seed) & 0.035 & 0.045 & 0.000 & 0.005 & \textbf{0.021} \\
\bottomrule
\end{tabular}      % generated by paper/scripts/make_figures.py
    }{%
        \fbox{\parbox{0.6\linewidth}{\centering\vspace{1cm}\textit{%
            tab1\_avg\_perf.tex placeholder -- copy from paper/}\vspace{1cm}}}%
    }
\end{table}

The headline finding (\Cref{tab:avg_perf}) is that the
unregularised Fine-Tuning control achieves an average success
rate of $0.139$, while every regularisation-based method reaches
at most $0.040$ --- a factor of three to seven worse. The
ranking among regularised methods (L2 ahead of MAS ahead of EWC,
within the precision the seed budget supports) matches the
ranking reported in the COOM paper, but the \emph{absolute}
position of the regularisation family relative to the
unregularised control inverts. Inversion, not re-ordering, is the
contribution. This addresses
requirement~\ref{req:multimethod}.

A useful sanity check on the magnitude of the gap is the EWC seed
band: across seeds 1, 2 and 3, EWC's CO4 average remains within
$\pm 0.012$ of its three-seed mean. The Fine-Tuning--vs.--EWC gap
of roughly $0.10$ is therefore at least eight times the
inter-seed spread of EWC. In other words, even allowing
generously for seed variance in the single-seed Fine-Tuning,
L2 and MAS runs, the qualitative ordering would not flip.

\subsection{Per-Task Curves}

\begin{figure}[t]
    \centering
    \includegraphics[width=\linewidth]{fig_success_curves}
    \caption{Per-task training success across the CO4 sequence,
        plotted within each task's own training window. The
        $y$-axis is the success rate; the $x$-axis is the epoch
        within the task ($1{,}000$ environment steps per epoch).
        EWC is plotted as a three-seed mean with a shaded band of
        $\pm$one standard deviation; the remaining methods are
        single-seed. Fine-Tuning matches the regularised methods
        on the first three tasks and dominates on Health
        Gathering, the last task in the sequence.}
    \label{fig:success_curves}
\end{figure}

\Cref{fig:success_curves} localises the gap to the final task in
the sequence, where Fine-Tuning reaches a last-five-epoch mean
success of $0.386$ while every regularised method stays at or
below $0.063$ (\Cref{tab:avg_perf}). On the first three tasks
the four methods are within one standard deviation of each
other, so the eventual gap in average performance is not caused
by an early-task disadvantage --- it is accumulated entirely on
the last task.

\subsection{Episode Return}

\begin{figure}[t]
    \centering
    \includegraphics[width=\linewidth]{fig_return_curves}
    \caption{Per-task training-time episode return for the four
        methods, plotted on the same axes as
        \Cref{fig:success_curves}. Return is the SAC reward signal
        the agent is actually optimising; success rate is the
        binary task-completion indicator derived from it. Where
        success rate stays at zero (Run-and-Gun for all methods)
        return is nonetheless non-trivial, indicating the policy
        is engaging with the environment but not crossing the
        success threshold.}
    \label{fig:return_curves}
\end{figure}

\Cref{fig:return_curves} reports the same runs measured by
episode return rather than task-success indicator. Two
observations follow. First, on Run-and-Gun all four methods
return between $25$ and $80$ reward per episode despite a zero
success rate; the policies are engaging with the environment
but not enough to clear the kill-count threshold that defines
success. The downstream domain-metric table
(\Cref{tab:domain_metrics}) confirms this by reporting a
non-zero kill count for each method on the same task. Second,
on Health Gathering the Fine-Tuning return rises monotonically
during the task while the regularised methods plateau within
the first ten epochs, mirroring the success-rate curves and
ruling out an explanation that the regularised agents are
optimising a different reward.

\subsection{Domain-Specific Per-Task Metrics}

\begin{table}[t]
    \centering
    \caption{Last-five-epoch mean of the per-task domain-specific
        metric, alongside the success rate already given in
        \Cref{tab:avg_perf}. Combat tasks (Chainsaw,
        Run-and-Gun) are scored by kills; survival-style tasks
        (Raise-the-Roof, Health Gathering) by the number of
        frames the agent stays alive. The Run-and-Gun column is
        the most informative: every method registers $0.8$--$2.2$
        kills despite zero success rate.}
    \label{tab:domain_metrics}
    \IfFileExists{tab5_domain_metrics.tex}{%
        \begin{tabular}{lcccc}
\toprule
Method & Chainsaw (kills) & Raise-the-Roof (frames) & Run-and-Gun (kills) & Health Gathering (frames) \\
\midrule
FT (1) & 1.9 & 205.6 & 2.2 & 338.5 \\
L2 (1) & 1.9 & 176.2 & 0.8 & 137.7 \\
MAS (1) & 0.0 & 202.9 & 1.5 & 135.9 \\
EWC (3) & 1.5 & 185.9 & 1.4 & 173.1 \\
EWC-Online (1) & 1.9 & 181.1 & 1.1 & 126.1 \\
\bottomrule
\end{tabular}%
    }{%
        \fbox{\parbox{0.7\linewidth}{\centering\vspace{1cm}%
            \textit{tab5\_domain\_metrics.tex placeholder.}%
            \vspace{1cm}}}%
    }
\end{table}

\Cref{tab:domain_metrics} adds a finer-grained reading of the
results, and the comparison to \Cref{tab:avg_perf} requires one
piece of context. The COOM benchmark normalises every per-task
success rate to the interval $[0,1]$ via the formula
\begin{equation}
    \mathrm{success}
    = \mathrm{clip}\!\left(
        \frac{m - m_{\mathrm{rand}}}{m_{\max} - m_{\mathrm{rand}}},
        \, 0, 1
    \right),
    \label{eq:succnorm}
\end{equation}
where $m$ is the per-task metric reported in
\Cref{tab:domain_metrics}, $m_{\mathrm{rand}}$ is the
expected value of the metric under a random action policy,
and $m_{\max}$ is the scenario's upper bound. For the four
CO4 tasks the relevant $(m_{\mathrm{rand}}, m_{\max})$ pairs
are $(10, 40)$ on Chainsaw, $(640, 2500)$ on Raise-the-Roof,
$(7.5, 20)$ on Run-and-Gun and $(800, 2500)$ on Health
Gathering. The clipping rule means that any agent performing
\emph{below} the random baseline records a success rate of
exactly zero, regardless of how close to the baseline it is.

Three observations follow directly. First, the column for
Run-and-Gun in \Cref{tab:avg_perf} is uniformly zero or
near-zero across all four methods, but
\Cref{tab:domain_metrics} shows that every method is killing
$0.8$--$2.2$ enemies per episode --- well \emph{below} the
$7.5$ kills the random policy achieves. The agents are not
inert; they are simply worse than random shooting on
Run-and-Gun at this budget, and \Cref{eq:succnorm} clips that
worse-than-random performance to a flat zero. Second, on
Chainsaw, MAS records zero kills on \Cref{tab:domain_metrics}
yet a non-zero last-five-epoch success ($0.057$ in
\Cref{tab:avg_perf}); this is consistent with a small number
of episodes within the averaging window having exceeded the
$10$-kill random baseline while the mean stays below it.
Third, on Health Gathering, all methods report mean episode
lengths between $135$ and $338$ frames, well below the
$800$-frame random baseline; Fine-Tuning's higher last-five
mean of $0.386$ in \Cref{tab:avg_perf} again reflects a
right-skewed distribution of per-episode lengths within the
averaging window rather than a uniform improvement.

\subsection{Learning Speed}

\begin{table}[t]
    \centering
    \caption{Number of in-task epochs needed to reach half of the
        in-task peak success rate. Lower is faster. ``---'' means
        the in-task peak success rate was zero.}
    \label{tab:learning_speed}
    \IfFileExists{tab3_learning_speed.tex}{%
        \begin{tabular}{lcccc}
\toprule
Method & T0 & T1 & T2 & T3 \\
\midrule
FT (1) & 18.0 & 14.0 & 19.0 & 16.0 \\
L2 (1) & 18.0 & 3.0 & --- & 3.0 \\
MAS (1) & 22.0 & 4.0 & --- & 10.0 \\
EWC (3) & 17.7 & 3.3 & 4.0 & 14.0 \\
EWC-Online (1) & 18.0 & 3.0 & --- & 3.0 \\
\bottomrule
\end{tabular}%
    }{%
        \fbox{\parbox{0.6\linewidth}{\centering\vspace{1cm}%
            \textit{tab3\_learning\_speed.tex placeholder.}%
            \vspace{1cm}}}%
    }
\end{table}

\Cref{tab:learning_speed} disentangles two aspects of
learning that the success-rate average conflates: \emph{how
quickly} a method climbs and \emph{how high} it gets. Fine-Tuning
takes $14$--$19$ in-task epochs to reach half its in-task peak;
the regularised methods reach half-peak in $3$--$10$ epochs on
most tasks, but their peak is much lower. In other words, the
regularised methods do not learn slowly --- they learn quickly
to a low ceiling and then stop. This pattern is consistent with
a regularisation penalty that suppresses gradient updates after
a brief initial phase of policy improvement, and it strengthens
the diagnosis in \Cref{sec:fishernoise}: it is not that EWC, L2
or MAS are too slow to converge in $20{,}000$ steps; it is that
they reach a fixed point that is below the Fine-Tuning peak and
stay there.

\subsection{Within-Task Drift}

\begin{table}[t]
    \centering
    \caption{Within-task drift, defined as the peak training
        success during a task minus the mean success in the last
        five epochs of the same task. A small value means the
        agent ends the task close to where it peaked; a large
        positive value means the agent regressed within the task.
        Computed from the same logs as \Cref{tab:avg_perf}.}
    \label{tab:forgetting}
    \IfFileExists{tab2_forgetting.tex}{%
        \begin{tabular}{lcccc|c}
\toprule
Method & T0 & T1 & T2 & T3 & \textbf{Avg.} \\
\midrule
FT (1) & 0.045 & 0.105 & 0.075 & 0.151 & \textbf{0.094} \\
L2 (1) & 0.049 & 0.063 & 0.000 & 0.060 & \textbf{0.043} \\
MAS (1) & 0.118 & 0.056 & 0.000 & 0.240 & \textbf{0.104} \\
EWC (3) & 0.072 & 0.034 & 0.004 & 0.171 & \textbf{0.070} \\
EWC-Online (1) & 0.049 & 0.053 & 0.000 & 0.066 & \textbf{0.042} \\
\bottomrule
\end{tabular}%
    }{%
        \fbox{\parbox{0.6\linewidth}{\centering\vspace{1cm}\textit{%
            tab2\_forgetting.tex placeholder -- run
            paper/scripts/make\_figures.py and re-copy.}%
            \vspace{1cm}}}%
    }
\end{table}

The drift values in \Cref{tab:forgetting} are mostly small
relative to the per-task success rates in \Cref{tab:avg_perf}.
The largest single drift, MAS on Health Gathering ($0.240$),
comes from an isolated success spike at one epoch within an
otherwise low-success task and is consistent with seed-level
exploration noise rather than systematic mid-task unlearning.
On the regularised methods that show non-trivial peak
performance, the drift is bounded by the success rate itself,
which means the gap in \Cref{tab:avg_perf} cannot be attributed
to within-task forgetting --- it is driven by the methods
failing to learn each task in the first place.

\subsection{Seed Variance for EWC}
\label{sec:seedband}

\begin{figure}[t]
    \centering
    \includegraphics[width=0.85\linewidth]{fig_ewc_seed_band}
    \caption{Last-five-epoch mean training success of EWC on each
        CO4 task, broken out by seed and compared to the
        three-seed mean. The inter-seed spread is at most
        ${\pm}0.033$ on Health Gathering and within ${\pm}0.012$
        on the other tasks.}
    \label{fig:ewc_seed_band}
\end{figure}

\Cref{fig:ewc_seed_band} reports per-seed last-five-epoch
success for the three EWC runs. The largest inter-seed spread
across the four tasks is ${\pm}0.033$ on Health Gathering, the
final task. The CO4 average across the three seeds is
$0.037 {\pm} 0.012$. The Fine-Tuning--vs.--EWC gap of
$0.139 - 0.037 = 0.102$ is therefore approximately eight times
the EWC inter-seed spread, which is the basis for the
robustness argument made earlier in this section.

\subsection{Regularisation Dynamics}

\begin{figure}[t]
    \centering
    \includegraphics[width=0.95\linewidth]{fig_lossreg}
    \caption{Regularisation loss
        $\mathcal{L}_{\mathrm{reg}}$ across the CO4 sequence,
        for the three regularised methods. The $x$-axis is the
        training epoch concatenated across the four tasks; the
        $y$-axis is on a symlog scale to make both the
        near-zero EWC values on the first task and the
        large-magnitude L2 values on later tasks visible in a
        single panel. Vertical dashed lines mark task
        boundaries. The penalty is exactly zero during the
        first task (no anchor exists yet), spikes at every
        subsequent task boundary, and grows monotonically
        across the sequence as additional importance terms are
        accumulated.}
    \label{fig:lossreg}
\end{figure}

\Cref{fig:lossreg} confirms that the regularisation mechanisms
operate as designed: $\mathcal{L}_{\mathrm{reg}}$ is zero
during the first task because no source-task importance
estimate exists yet, jumps at every subsequent task boundary as
the new anchor activates, and grows monotonically across the
sequence. The failure mode is therefore behavioural rather than
implementation-level; this addresses
requirement~\ref{req:diagnose}.

\subsection{Diagnosis: Fisher on Noise}
\label{sec:fishernoise}

The remaining question is \emph{why} a correctly firing
regularisation penalty depresses average success below the
no-regularisation control. The hypothesis adopted here is the
``Fisher on noise'' effect: at $20{,}000$ steps per task, the
source task is not converged, so the parameter vector at the
end of the task is a point on a noisy optimisation trajectory
rather than a local optimum. The Fisher diagonal evaluated at
that point therefore captures the curvature of the loss in the
direction of recent gradient noise rather than the curvature
around a learned solution. The same mechanism applies to
MAS, whose importance signal is the gradient norm of the model
output at end-of-task. When the agent's parameters are
subsequently anchored to that noisy point on the next task, the
anchor is to a noise direction rather than to a meaningful
solution, and the effective constraint is to suppress
gradient updates without a corresponding stability benefit.

Three observations support this diagnosis. First, the
within-task drift in \Cref{tab:forgetting} is small, which
rules out mid-task unlearning as the explanation: regularised
methods are not failing because they are forgetting the new
task, they are failing because they never learn it well in the
first place. Second, the gap concentrates on the final task
(Health Gathering, \Cref{fig:success_curves}), which is the
task that has accumulated the most prior anchors and therefore
the most exposure to the cumulative noise. Third, the
regularisation losses themselves
(\Cref{fig:lossreg}) grow super-linearly across the sequence,
consistent with monotonic accumulation of Fisher matrices that
each contribute a noise floor.

\begin{figure}[t]
    \centering
    \includegraphics[width=\linewidth]{fig_training_dynamics}
    \caption{SAC internals across the CO4 sequence, comparing
        Fine-Tuning to EWC. Top: actor loss $\mathcal{L}_\pi$;
        middle: $Q_1$ loss $\mathcal{L}_{Q_1}$; bottom: SAC's
        auto-tuned entropy temperature $\alpha$ for the
        currently active task. Vertical dashed lines mark task
        boundaries. The plot is intended as evidence that the
        source-task end state on which Fisher is evaluated is
        far from converged: $\mathcal{L}_{Q_1}$ continues to
        rise across the sequence rather than plateauing, and
        $\alpha$ is reset to its initial value at every task
        switch and only partially decays before the next switch.}
    \label{fig:training_dynamics}
\end{figure}

\Cref{fig:training_dynamics} adds a fourth, more direct piece
of evidence. The plot tracks SAC's actor loss, $Q_1$ critic
loss and entropy temperature $\alpha$ for Fine-Tuning and EWC
across the full sequence. Two features stand out. The $Q_1$
loss rises monotonically across the sequence rather than
plateauing within each task, which implies that the value
estimate at the end of every task is still drifting --- exactly
the regime in which the Fisher estimate is least
representative of a stable optimum. The entropy temperature
$\alpha$, which SAC auto-tunes per task, is reset to its
initial value $e \approx 2.72$ at every task boundary and only
partially decays before the next switch. Anchoring on a
parameter vector taken at this transient state is precisely the
``Fisher on noise'' regime the diagnosis predicts.

\begin{figure}[t]
    \centering
    \includegraphics[width=\linewidth]{fig_action_dist}
    \caption{Action-distribution share, averaged across the
        last five epochs of each task per method. The 12 raw
        ViZDoom actions are grouped into NO-OP/EXECUTE,
        MOVE\_FORWARD-and-its-variants,
        TURN\_RIGHT-and-its-variants, and
        TURN\_LEFT-and-its-variants. Bars sum to one within
        each method-task pair.}
    \label{fig:action_dist}
\end{figure}

\Cref{fig:action_dist} provides a complementary view of what
the policies actually \emph{do}. None of the four methods
collapses to a degenerate single-action policy on any task: the
action distribution is non-uniform but never one-hot, and the
share of MOVE\_FORWARD on Health Gathering grows for
Fine-Tuning relative to the regularised methods, consistent
with Fine-Tuning's higher episode-length and success rate on
that task. This rules out catastrophic policy collapse as an
alternative explanation for the gap reported in
\Cref{tab:avg_perf}.

\begin{figure}[t]
    \centering
    \includegraphics[width=\linewidth]{fig_action_entropy}
    \caption{Shannon entropy of the 12-action policy
        distribution per epoch, across the full CO4 sequence.
        The dotted horizontal line marks the entropy of a
        uniform 12-action policy ($\log_2 12 \approx 3.58$);
        any policy at or near that line is essentially
        random. Vertical dashed lines mark task boundaries.
        Both Fine-Tuning and EWC stay well below the uniform
        line throughout, confirming that the policy is
        committing to specific actions rather than guessing,
        but stays well above zero, confirming that the policy
        is not collapsing to a single deterministic action
        either. The entropy spikes at every task boundary as
        the per-task SAC entropy temperature $\alpha$ is
        re-tuned, then decays.}
    \label{fig:action_entropy}
\end{figure}

\Cref{fig:action_entropy} extends the action-share view of
\Cref{fig:action_dist} into a per-epoch trajectory: the
Shannon entropy of the action distribution drops smoothly
within each task as the policy specialises, then jumps back
toward uniform at every task boundary as SAC's entropy
temperature is re-tuned. The decay is monotonically less
deep over successive tasks for EWC than for Fine-Tuning,
which is consistent with EWC's regularisation penalty
suppressing the gradient updates that drive policy
specialisation.

The diagnosis is testable: at the COOM paper's $200{,}000$-step
budget the source task is closer to convergence and the Fisher
estimate should be cleaner. This is the future work item in
\Cref{sec:future}.

\section{Improved EWC Variant: Online EWC}
\label{sec:ewc_online}

The Fisher-on-noise diagnosis above motivates a single,
minimally invasive change to EWC. \Cref{sec:ewc_online_impl}
documents the implementation: instead of accumulating per-task
Fisher diagonals additively as in \cref{eq:fisher_sum}, the
running estimate is updated by the exponential moving average
of \cref{eq:fisher_ema}, with decay factor $\gamma = 0.95$
inherited from the original Online EWC proposal of
\citet{schwarz2018progress}. All other hyperparameters --- the
regularisation coefficient $\lambda = 250$, the SAC learning
rate, the replay buffer reset behaviour --- are held fixed at
the values used by the original EWC.

\begin{table}[t]
    \centering
    \caption{Online EWC vs original EWC at the same CO4 /
        $20{,}000$-step budget. Last-five-epoch mean training
        success per task. Values for original EWC are
        reproduced from \Cref{tab:avg_perf} for ease of
        comparison; the Online EWC row is added in this
        section.}
    \label{tab:ewc_online}
    \IfFileExists{tab6_ewc_online.tex}{%
        \begin{tabular}{lcccc|c}
\toprule
Method & T0 & T1 & T2 & T3 & \textbf{Avg.\ Perf.} \\
\midrule
FT (1 seed) & 0.054 & 0.098 & 0.019 & 0.386 & \textbf{0.139} \\
EWC (3 seeds) & 0.028 $\pm$ 0.012 & 0.056 $\pm$ 0.005 & 0.000 $\pm$ 0.000 & 0.063 $\pm$ 0.033 & \textbf{0.037} \\
EWC-Online (1 seed) & 0.035 & 0.045 & 0.000 & 0.005 & \textbf{0.021} \\
\bottomrule
\end{tabular}%
    }{%
        \fbox{\parbox{0.7\linewidth}{\centering\vspace{1cm}%
            \textit{tab6\_ewc\_online.tex placeholder ---
            generated automatically once the Online EWC
            CO4/20k seed-1 run completes
            (\texttt{logs/ewc\_online\_co4\_20k\_seed1/}).}%
            \vspace{1cm}}}%
    }
\end{table}

\begin{figure}[t]
    \centering
    \IfFileExists{figures/fig_ewc_vs_online.pdf}{%
        \includegraphics[width=\linewidth]{fig_ewc_vs_online}%
    }{%
        \fbox{\parbox{\linewidth}{\centering\vspace{2cm}%
            \textit{fig\_ewc\_vs\_online placeholder ---
            generated automatically once the Online EWC run
            completes.}\vspace{2cm}}}%
    }
    \caption{Per-task training success curves for original EWC
        (mean across three seeds, shaded band $\pm$ one
        standard deviation) overlaid against Online EWC
        (single seed). The two methods are run with identical
        hyperparameters except the Fisher-merge rule.}
    \label{fig:ewc_vs_online}
\end{figure}

The expected behaviour, predicted by the diagnosis in
\Cref{sec:fishernoise}, is that the EMA bounds the cumulative
regularisation tax across tasks and that the noise component
of any single Fisher estimate contributes only $1 - \gamma$ of
its magnitude, so Online EWC should improve over the original
EWC at this budget by a small but visible margin. The
quantitative reading of \Cref{tab:ewc_online} below is filled
in once the Online EWC run completes; the discussion that
follows here covers the three possible outcomes and the
narrative each supports.

\paragraph{Outcome A: Online EWC improves on EWC.} If the
Online EWC last-five-epoch CO4 average exceeds the
$0.037 {\pm} 0.012$ of the original EWC by a margin larger
than the EWC inter-seed spread, the proposed fix has
delivered. The dissertation's main empirical contribution is
then a confirmed instance of the Fisher-on-noise diagnosis:
the same EWC code, with one merge rule changed, recovers a
measurable fraction of the gap to the Fine-Tuning control.

\paragraph{Outcome B: Online EWC matches EWC within the seed
band.} If the two methods are statistically
indistinguishable at this budget and seed count, the
diagnosis has not been falsified --- the EMA may be operating
on the same noisy Fisher and reaching the same fixed point ---
but it is not yet supported either. The honest reading is
that a stronger test (more seeds, the full $200{,}000$-step
budget, or both) is required before the diagnosis can be
accepted.

\paragraph{Outcome C: Online EWC underperforms EWC.} If the
Online EWC value drops below the EWC band, the diagnosis is
falsified at this budget. Two readings are possible: the EMA
is too aggressive (a smaller $\gamma$ closer to $1$ would
restore performance) or the original additive Fisher was, in
fact, accidentally the right kind of noise floor for low-budget
training. Either reading is publishable as a negative result.

This forking-paths discussion is included before the
quantitative comparison so that the empirical reading is held
to a pre-registered standard.

\subsection{Paper-Budget Cross-Check}
\label{sec:paper_budget}

The CPU runs are conducted at $1/10$ the per-task budget of the
COOM paper. To check that the gap to the paper's reported
EWC score of $0.54$ is genuinely budget-driven (rather than an
artefact of the CPU build, the missing held-out evaluation, or
the seed selection), single-seed GPU runs of EWC \emph{and the
proposed Online EWC variant} at the published $200{,}000$-step
budget were executed on a rented RTX 3090, this time with the
held-out evaluation flag enabled.
\Cref{tab:paper_budget} reports the held-out CO4 success per
task and the four-task average for both methods.

\begin{table}[t]
    \centering
    \caption{Paper-budget reproduction: a single seed of EWC
        run at the COOM paper's published $200{,}000$
        steps-per-task on the CO4 sequence, with held-out
        evaluation enabled. Training-time numbers are
        comparable to \Cref{tab:avg_perf}; held-out numbers
        are the metric the COOM paper itself reports.}
    \label{tab:paper_budget}
    \IfFileExists{tab7_paper_budget.tex}{%
        \begin{tabular}{lcccc|c}
\toprule
Method & T0 & T1 & T2 & T3 & \textbf{Avg.} \\
\midrule
EWC & 0.993 & 0.000 & 0.045 & 0.011 & \textbf{0.262} \\
EWC-Online ($\gamma$=0.50) & 0.998 & 0.042 & 0.000 & 0.019 & \textbf{0.265} \\
\bottomrule
\end{tabular}%
    }{%
        \fbox{\parbox{0.7\linewidth}{\centering\vspace{1cm}%
            \textit{tab7\_paper\_budget.tex placeholder.}%
            \vspace{1cm}}}%
    }
\end{table}

\begin{figure}[t]
    \centering
    \includegraphics[width=\linewidth]{fig_paper_budget}
    \caption{Per-task held-out (blue solid) and
        training-time (grey dashed) success rates across the
        full paper-budget EWC sequence. Vertical dotted
        lines mark approximate task boundaries. The held-out
        line on Chainsaw stays close to $1.0$ throughout the
        sequence, the held-out line on Raise-the-Roof
        collapses to zero immediately after that task ends,
        and the held-out lines on the two later tasks track
        their training-time signal closely.}
    \label{fig:paper_budget}
\end{figure}

Two readings of \Cref{tab:paper_budget} matter for the
dissertation.

\paragraph{Pipeline-level reproduction is strengthened at
paper budget.} The CO4 average for the original EWC climbs
from $0.037$ at the $20{,}000$-step CPU budget to a held-out
$0.262$ at $200{,}000$ steps --- a roughly seven-fold
recovery without any method change. Both numbers remain
below the COOM paper's reported $0.54$, but the gap is
consistent with the differences this study still does not
control: a single seed against the paper's ten, a single
sequence (CO4) against the paper's seven, and a per-task
budget that is the paper's smallest sequence average rather
than its full per-task allocation. The principal qualitative
observation is that paper-budget EWC behaves like the COOM
paper's EWC and not like our CPU EWC: Chainsaw is retained
almost perfectly ($0.993$ held-out at the end of the
sequence), Raise-the-Roof is catastrophically forgotten
($0.000$), and the later tasks sit between the two. The
pipeline-level reproduction claim of
\Cref{sec:reproduction_scope} is therefore strengthened.

\paragraph{Online EWC at $\gamma = 0.5$ flips from
Outcome~C to Outcome~A when the budget moves from
$20{,}000$ to $200{,}000$.} At the CPU budget, Online EWC
with $\gamma = 0.95$ underperformed the original EWC
(\Cref{sec:ewc_online_outcome}, training-time average
$0.021$ vs.\ $0.037$). At the paper budget reported here,
Online EWC with $\gamma = 0.5$ \emph{slightly outperforms}
the original EWC on the held-out CO4 average ($0.265$ vs.\
$0.262$), with the per-task improvement concentrated on
Raise-the-Roof: catastrophic forgetting ($0.000$) is
replaced by a small positive retention ($0.042$) under the
EMA merge rule. The improvement is within the seed
variance recorded for the original EWC at the CPU budget
($\pm 0.012$), so a multi-seed extension is needed before
the gap can be called statistically separable. What can be
said now, however, is that the qualitative ordering between
EWC and Online EWC predicted by the Fisher-on-noise
diagnosis (\Cref{sec:fishernoise}) does emerge once the
source task is given enough budget to converge, exactly as
the diagnosis claimed.

\subsection{Observed Outcome}
\label{sec:ewc_online_outcome}

The single-seed Online EWC run completed cleanly on the same
CO4 / $20{,}000$-step configuration as the original EWC. The
last-five-epoch CO4 average is reported in
\Cref{tab:ewc_online}: Online EWC reaches $0.021$, against
the original EWC's three-seed $0.037 {\pm} 0.012$. The
Online EWC value falls below the lower edge of the EWC seed
band ($0.025$), placing this result in
\textbf{Outcome~C} of the pre-registered analysis above ---
the EMA variant \emph{underperforms} the additive baseline
at $\gamma = 0.95$. The gap is concentrated on the final
task (Health Gathering, $0.005$ for Online EWC vs.\
$0.063 {\pm} 0.033$ for original EWC), with the early
tasks essentially tied and Run-and-Gun clipped to zero for
both methods by the COOM normalisation
(\Cref{sec:reproduction_scope}, \cref{eq:succnorm}).

The honest reading is that the diagnosis was \emph{not}
falsified --- in particular, the gap concentrates on the
final task, which is where the diagnosis predicted the
constraint dynamics matter most --- but the specific
intervention chosen ($\gamma = 0.95$) overshoots the right
operating point at this budget and seed count. Two readings
of the result are consistent with the underlying mechanism:
\emph{(a)} the EMA at $\gamma = 0.95$ preserves an
unusually large weight on the very first Fisher estimate
($\gamma^{2} \approx 0.90$ even three task transitions
later), and that first Fisher is precisely the noisiest one
in the diagnosis, so the EMA is amplifying the noise it was
meant to attenuate; or \emph{(b)} the single-seed Online EWC
run landed in the lower tail of its seed distribution on
Health Gathering, where the original EWC band already shows
${\pm}0.033$ inter-seed spread.

A $\gamma$ sweep at this budget --- in particular,
$\gamma \in \{0.5, 0.7\}$, which weight the first Fisher
estimate at $0.25$ and $0.49$ respectively rather than at
$0.90$ --- and a multi-seed extension at the chosen
$\gamma$ are the two follow-on items required before either
reading can be selected. Both are explicitly listed in
\Cref{sec:future}.

\paragraph{A bonus result.}
A small unanticipated finding emerges from the same run:
Online EWC completed the CO4 sequence in $48$~minutes,
against $128 {\pm} 8$~minutes for the original EWC
(\Cref{tab:walltime}). The $2.7{\times}$ speedup arises
because the EMA merge replaces an unbounded-magnitude
additive accumulation with a rescale-and-add, which
keeps the regularisation-loss arithmetic in a
small-magnitude regime throughout the sequence. The
finding is incidental to the dissertation's main argument
but is recorded here because it bears on the original
\Cref{tab:walltime} discussion of compute-aware method
design.

\section{Threats to Validity}
\label{sec:validity}

The study is deliberately scoped to a single sequence, a single
budget, and a single hardware platform. Each is a threat to the
generality of the conclusion.

\paragraph{Single sequence.}
CO4 is one of seven sequence families in COOM. The published
ranking is averaged across all seven; restricting to CO4 is a
real loss of generality, although CO4 is also the sequence that
the published methods score \emph{best} on relative to
Fine-Tuning, so if anything it is biased \emph{against} the
finding reported here. A study at the published $200{,}000$
budget but restricted to CO4 would resolve this.

\paragraph{Single seed for Fine-Tuning, L2 and MAS.}
The seed band is anchored only on EWC. The argument that the
ranking does not flip even allowing generously for seed variance
relies on the assumption that the inter-seed spread of EWC
(${\pm}0.012$) is representative of the spread for the other
methods. This assumption is plausible because all four methods
share the same SAC backbone and replay buffer reset behaviour,
but it is an assumption rather than a measurement.

\paragraph{Training-time signal rather than held-out
evaluation.}
The COOM paper's Forgetting and Forward Transfer metrics
require evaluating the agent on past tasks during training, via
the upstream \texttt{--test} flag. This study runs with
\texttt{--no\_test} to keep the overnight queue under three
hours per run, and reports last-five-epoch training success
instead. Held-out success is generally lower than training
success because no exploration noise is injected at evaluation
time, but the relative ordering between methods is preserved in
the COOM paper's own data, which is the property used here.

\paragraph{Single hardware platform.}
All runs are on a single Apple-Silicon CPU host. CPU floating
point on this platform differs from the CUDA numerics under
which the COOM hyperparameters were originally tuned. A small
quantitative difference in absolute success rates between this
study and the original is therefore expected; the magnitude of
the inter-method gap reported in \Cref{tab:avg_perf} is too
large to be explained by this difference alone.

\paragraph{Hyperparameters not retuned.}
The regularisation strengths $\lambda_{\text{EWC}}=250$,
$\lambda_{\text{L2}}=10^{5}$ and $\lambda_{\text{MAS}}=10^{4}$
are inherited verbatim from the COOM paper rather than retuned
at the new budget. It is therefore possible that a smaller
$\lambda$ tuned for $20{,}000$ steps per task would close the
gap to Fine-Tuning. The same applies to the Online EWC decay
factor $\gamma = 0.95$ (\Cref{sec:ewc_online}), which was
inherited from \citet{schwarz2018progress} where it was used
for a within-task EMA refresh schedule rather than the
per-task merge schedule used here. A $\gamma$ sweep at this
budget is the third item in \Cref{sec:future} and is exactly
the ``compute-aware regularisation schedule'' referenced in
the abstract.

\paragraph{Online EWC reported as a single seed.}
Within the project window, Online EWC was run on a single
seed (seed~$1$) on the CO4 / $20{,}000$-step configuration.
The original EWC reports a three-seed band of ${\pm}0.012$ on
the CO4 average (\Cref{fig:ewc_seed_band}); applying that
inter-seed spread to the single Online EWC seed gives a
robustness margin against which the head-to-head should be
read. The full three-seed run for Online EWC is the first
item in \Cref{sec:future} and is required before any
publication-grade claim about the method's effect can be
made. The forking-paths analysis in \Cref{sec:ewc_online} is
written so that each of the three possible outcomes carries
the appropriate hedge.
